\hypertarget{ii.-heffingen}{%
\section{II. Heffingen}\label{ii.-heffingen}}

\hypertarget{iia.-administratieve-heffingen}{%
\subsection{IIa. Administratieve
heffingen}\label{iia.-administratieve-heffingen}}

\hypertarget{art.-2}{%
\paragraph{Art. 2}\label{art.-2}}

Een administratieve heffing wordt geheven op elke geldelijke straf of
schadevergoeding uitgesproken door de Linker.

\hypertarget{art.-3}{%
\paragraph{Art. 3}\label{art.-3}}

Een administratieve heffing mag ten hoogste 10\% zijn van het
verschuldigde bedrag.

\hypertarget{art.-4}{%
\paragraph{Art. 4}\label{art.-4}}

Indien een administratieve heffing nog niet geint is voor het einde van
de parlementaire zitting zal deze eveneens gebruikt worden om:

\begin{itemize}
\tightlist
\item
  \emph{In het geval dat deze is afgehouden van een geldboete}: mee te
  nuttigen tijdens het nuttigen van de Eurodeals van dat jaar.
\item
  \emph{In het geval dat deze is afgehouden van een schadevergoeding}:
  uit betaald to worden aan de bevoordeelde van de schadevergoeding.
\end{itemize}

\hypertarget{iib.-vrijwillige-heffingen}{%
\subsection{IIb. Vrijwillige
heffingen}\label{iib.-vrijwillige-heffingen}}

\hypertarget{art.-5}{%
\paragraph{Art. 5}\label{art.-5}}

Indien het Ministerie van Alles ambtshalve een aankoop moet doen voor de
Vrijstaat kan deze, indien de aankoop niet volledig gedekt is door de
nog openstaande administratieve heffingen, een Aanvraag tot intekening
op een Vrijwillige Heffing uitschrijven.

\hypertarget{art.-5bis}{%
\paragraph{Art. 5bis}\label{art.-5bis}}

Een Aanvraag tot intekening op een Vrijwillige Heffing zal ten minste
bestaan uit:

\begin{itemize}
\tightlist
\item
  Het overblijvende aankoopbedrag na aftrek van de nog openstaande
  administratieve heffingen.
\item
  Een korte beschrijving van de aankoop.
\item
  De verdeelsleutel tussen de inschrijvers, standaard zal deze (1/\#
  inschrijvers) zijn
\item
  Indien het aankoopbedrag boven de 50 Eurodeals wordt geschat, een
  vergelijking tussen ten minste drie leveranciers. De beste leverancier
  wordt gekozen door een objectief punten systeem. Standaard zal dit
  gebaseerd zijn op de laagste kostprijs.
\end{itemize}

\hypertarget{art.-6}{%
\paragraph{Art. 6}\label{art.-6}}

Elke Idioot zal de Aanvraag tot intekening op een Vrijwillige Heffing
ontvangen via fax, briefpost, elektronische post of per postduif indien
beschikbaar.

\hypertarget{art.-7}{%
\paragraph{Art. 7}\label{art.-7}}

Elke Idioot zal na ontvangst van de Aanvraag tot intekening op een
Vrijwillige Heffing een antwoord sturen per fax, briefpost of
elektronische post met daarin:

\begin{itemize}
\tightlist
\item
  De volgnummer van de heffing
\item
  Het antwoord van de Idioot, dit kan \emph{ja} of \emph{nee} zijn.
\end{itemize}

Indien niet aan bovenstaand formaat is voldaan is het antwoord niet
geldig. Antwoorden op verschillende Aanvragen tot intekening op een
Vrijwillige Heffing mogen op hetzelfde blad geschreven worden zolang het
dossiernummer duidelijk is voor elke \emph{ja} of \emph{nee}.

Indien het Ministerie van Alles geen geldig antwoord heeft ontvangen van
een Idioot per fax, briefpost of elektronische post vier weken na het
versturen van de Aanvraag tot intekening op een Vrijwillige Heffing zal
de Idioot aanzien worden als ingeschreven op de heffing.
